\scanpage{31}
% letterhead
\noindent\begin{minipage}[t]{0.57\linewidth}
\centering\rmfamily
UNIVERSITY OF SCIENCE AND TECHNOLOGY\\
OF LANGUEDOC\\
\rule{0.55cm}{0.4pt}\\
\sffamily
INSTITUTE OF MATHEMATICS\\
Place Eugène Bataillon\\
34060 MONTPELLIER CEDEX\\
\rule{0.65cm}{0.4pt}
\end{minipage}\hfill
\begin{minipage}[t]{0.37\linewidth}
Montpellier, the
\end{minipage}

\begin{center}
\setlength{\ULdepth}{1.45\baselineskip}
\uline{COURSE\quad C4 1977-78 by $\mathbb A$. GROTHENDIECK}\par\vspace{.6\baselineskip}
\end{center}

\uline{Geometry of the icosahedron}

It is strange that the icosahedron, studied for more than two thousand years in very diverse contexts and from very diverse viewpoints, seems never to have been studied systematically from the combinatorial point of view, nor from that of geometric realizations over general fields (or even rings). This ``course'' would like to be a joint investigation of this theme, starting from the most naïve combinatorial intuition, exemplified on cardboard models (with a suitable colouring of the edges).

We shall therefore approach the theory of the i\deleted{e}cosahedron from a combinatorial point of view first, making explicit in detail the dictionary between the combinatorial geometry of the oriented icosahedron, and that of the ``oriented'' set with five elements (the five colors colouring the edges !) - which will at the same time provide a motivation and a geometric inspiration for studying the automorphism group common to the two situations, namely the alternating group%
% Detached pair at approximately (1582,1659) and (1577,1666), between this source line and the next; hand and verbal role unassigned.
\margin{\figure{\begin{tikzpicture}[x=1mm,y=1mm]
\fill (0.60,0.90) .. controls (0.37,0.85) and (0.39,0.59) .. (0.57,0.58) .. controls (0.77,0.59) and (0.86,0.82) .. (0.60,0.90) -- cycle;
\fill (0.00,0.00) .. controls (-0.25,0.03) and (-0.24,-0.16) .. (-0.06,-0.22) .. controls (0.13,-0.21) and (0.22,-0.07) .. (0.00,0.00) -- cycle;
\end{tikzpicture}}}%
 $\handwritten{\unsure{\mathfrak A|\mathcal A}}_5$ (which is also the smallest non-abelian simple group). The question of constructing ``geometric realizations'' of the combinatorial icosahedron, as a ``regular polyhedron'' over the field of real numbers, or over any other field, will be an opportunity both to deepen the understanding of the combinatorial structure, and to familiarize ourselves with the notions of vector and affine geometry (or even projective geometry) over fields of arbitrary characteristic.

% Detached dark shape at the lower-left physical edge area, approximately [20,2293,49,2317]; hand and physical origin unassigned.
\par\vspace{2\baselineskip}
\noindent\hspace*{-1.5cm}\begin{minipage}{0.8cm}
\figure{\begin{tikzpicture}[x=1mm,y=1mm]
\fill (0.00,0.00) .. controls (-0.16,0.17) and (-0.03,0.40) .. (0.13,0.48)
 -- (0.17,0.72) -- (0.41,0.69) -- (0.50,1.05) -- (0.76,1.19)
 -- (0.98,1.08) -- (1.20,1.35) -- (1.48,1.48) -- (1.64,1.72)
 -- (1.91,1.86) -- (2.16,1.84) -- (2.42,2.02) -- (2.47,1.85)
 -- (2.71,1.74) -- (2.78,1.28) -- (2.50,1.25) -- (2.36,1.02)
 -- (2.08,1.09) -- (1.90,0.84) -- (1.70,0.82) -- (1.61,0.43)
 -- (1.37,0.32) -- (1.25,0.00) -- (0.95,-0.08) -- (0.73,-0.24)
 -- (0.50,-0.10) .. controls (0.24,-0.13) and (0.07,-0.12) .. (0.00,0.00) -- cycle;
\end{tikzpicture}}
\end{minipage}
% archive sheet number: 30
